\documentclass[11pt]{article}
\usepackage[a4paper,margin=0.9in]{geometry}
\usepackage[T1]{fontenc}
\usepackage{textcomp}
\usepackage{helvet}
\renewcommand{\familydefault}{\sfdefault}
\usepackage{xcolor}
\usepackage{booktabs}
\usepackage{array}
\usepackage{colortbl}
\usepackage{float}
\usepackage{amsmath}
\usepackage{titlesec}
\usepackage{fancyhdr}
\usepackage{enumitem}
\usepackage{lastpage}
\usepackage[hidelinks]{hyperref}

\definecolor{bcnavy}{RGB}{11,31,58}
\definecolor{bcblue}{RGB}{32,74,135}
\definecolor{bcgold}{RGB}{176,141,87}
\definecolor{bcgrey}{RGB}{90,98,112}
\definecolor{bcrule}{RGB}{210,214,220}
\definecolor{bckeep}{RGB}{232,237,244}
\definecolor{bcact}{RGB}{247,241,228}

\titleformat{\section}{\color{bcnavy}\Large\bfseries}{\thesection}{0.6em}{}
\titleformat{\subsection}{\color{bcblue}\large\bfseries}{\thesubsection}{0.6em}{}
\titlespacing*{\section}{0pt}{1.3em}{0.45em}

\pagestyle{fancy}\fancyhf{}
\renewcommand{\headrulewidth}{0pt}
\renewcommand{\footrulewidth}{0.4pt}
\renewcommand{\footrule}{\hbox to\headwidth{\color{bcrule}\leaders\hrule height \footrulewidth\hfill}}
\fancyfoot[L]{\footnotesize\color{bcgrey}Bayesian Capital}
\fancyfoot[C]{\footnotesize\color{bcgrey}Confidential --- Internal}
\fancyfoot[R]{\footnotesize\color{bcgrey}Page \thepage\ of \pageref{LastPage}}

\setlength{\parindent}{0pt}\setlength{\parskip}{0.5em}

\newcommand{\note}[2]{%
  \begin{center}\fcolorbox{bcrule}{#1}{\parbox{0.95\textwidth}{\vspace{0.2em}#2\vspace{0.2em}}}\end{center}}

\begin{document}

{\color{bcgold}\rule{\textwidth}{1.6pt}}\\[0.3cm]
{\color{bcnavy}\huge\bfseries Trading Status}\\[0.15cm]
{\color{bcblue}\Large The Company's Activity Against the Badges of Trade}\\[0.35cm]
{\color{bcgold}\rule{\textwidth}{0.8pt}}\\[0.35cm]

\begin{tabular}{@{}ll@{}}
\textbf{Company} & Bayesian Capital Limited, company number 815246 \\
\textbf{Registered office} & 76 Georges Avenue, Blackrock, Co. Dublin \\
\textbf{Date} & 5 October 2026 \\
\textbf{Trading record to} & 2 October 2026 \\
\textbf{Source} & \texttt{TradingExcel\_9stock\_avgo.xlsx}, Trade Log and Discretionary Trade Log \\
\textbf{Status} & Internal record --- maintained monthly \\
\end{tabular}

\vspace{0.6em}
{\color{bcrule}\rule{\textwidth}{0.6pt}}

\section{Position}

\textbf{The company carries on a trade of dealing in listed shares.} Its sole activity is the
purchase and resale of fully paid ordinary shares listed on United States exchanges. It holds no
other assets and carries on no other activity.

Three facts define the operation, and each is a matter of record rather than of intention:

\begin{itemize}[leftmargin=1.4em,itemsep=0.2em]
\item \textbf{Every position's exit price is fixed before the position exists.} All 109 positions
      opened to date, across both the systematic and the discretionary books, had their selling
      price determined before the shares were acquired.
\item \textbf{Every purchase is made below the prevailing market price}, at a limit price computed
      in advance. The company does not buy at market and does not raise its bid to obtain stock.
\item \textbf{No position may be held beyond 50 calendar days.} The rule has been in force since
      June 2026 and no exception has been made to it.
\end{itemize}

In the nine weeks from 4 August to 2 October 2026 --- 44 trading days --- the company opened
\textbf{109 positions} across all nine shares in its universe and closed \textbf{97}. Thirty-seven
closed on the day they were bought and 77 within three days. The median completed holding period
was \textbf{one day}; the longest was 31. Two of 97 completed positions realised a loss. Gross
purchases were \textbf{USD 35,581,460} and realised profit \textbf{USD 780,702}, across nine
consecutive profitable weeks.

The purchase and the resale are a single decision, taken before the asset is acquired. The company
has never bought a share without having first fixed the price at which it intends to sell it.

\section{The trading record}

\begin{center}
\footnotesize
\begin{tabular}{@{}l r r r@{}}
\toprule
\textbf{4 Aug --- 2 Oct 2026} & \textbf{Systematic} & \textbf{Discretionary} & \textbf{Combined} \\
\midrule
Positions opened                  & 89 & 20 & \textbf{109} \\
Positions closed                  & 81 & 16 & \textbf{97} \\
Open at 2 October                 &  8 &  4 & 12 \\
Exit price fixed before entry     & 89 of 89 & 20 of 20 & \textbf{109 of 109} \\
Closed same day as bought         & 30 (37\%) & 7 (44\%) & 37 (38\%) \\
Closed within three days          & 64 (79\%) & 13 (81\%) & 77 (79\%) \\
Median completed hold             & 1 day & 1 day & \textbf{1 day} \\
Longest completed hold            & 31 days & 22 days & 31 days \\
Completed positions at a loss     & 2 & 0 & 2 of 97 \\
Distinct shares dealt in          & 9 & 4 & 9 of 9 \\
Gross value of purchases          & USD 27,699,177 & USD 7,882,282 & USD 35,581,460 \\
Largest single position           & USD 614,378 & USD 825,030 & USD 825,030 \\
Realised profit                   & USD 569,186 & USD 211,516 & \textbf{USD 780,702} \\
\bottomrule
\end{tabular}
\end{center}

Activity is continuous rather than episodic. The company placed purchases on 36 of the 44
trading days in the period, a rate of 2.48 new positions per trading day sustained since launch,
and dealt in every one of the nine shares in its universe.

\subsection{Realised profit by week}

\begin{center}
\footnotesize
\begin{tabular}{@{}l r r@{}}
\toprule
\textbf{Week ending} & \textbf{Realised P\&L (USD)} & \textbf{Trades} \\
\midrule
7 August 2026    &  59,691 &  9 \\
14 August 2026   &  25,747 &  8 \\
21 August 2026   & 145,807 & 13 \\
28 August 2026   &  89,554 &  8 \\
4 September 2026 & 109,452 &  9 \\
11 September 2026 &  65,284 &  8 \\
18 September 2026 &  43,389 &  6 \\
25 September 2026 &  73,176 & 11 \\
2 October 2026    & 168,602 &  9 \\
\midrule
\textbf{Total}    & \textbf{780,702} & \textbf{81} \\
\bottomrule
\end{tabular}
\end{center}

\subsection{Positions by share}

\begin{center}
\footnotesize
\begin{tabular}{@{}l r r r r r r r r r@{}}
\toprule
& VST & TSM & MU & GM & MRVL & VRT & VLO & AVGO & CF \\
\midrule
Positions opened & 19 & 19 & 14 & 13 & 13 & 10 & 9 & 8 & 4 \\
\bottomrule
\end{tabular}
\end{center}

\section{How the company operates}

\subsection{The exit is fixed before the entry}

This is the central feature of the operation and it applies without exception to both books.

\textbf{Systematic book.} Each morning two independent quantitative models compute, for each share,
a limit price below an estimated fair value. On execution of a purchase, a corresponding sell order
is placed at a price determined at the moment of entry, as a fixed percentage above the purchase
price. The sell order rests with the broker from the moment the position exists.

\textbf{Discretionary book.} The target sale price is recorded in the Discretionary Trade Log
before the buy order is transmitted. The log carries a dedicated \emph{Target sale} field and it
is populated on all 20 positions taken to date. The standard uplift is 5.0\% above the purchase
price; tighter targets are set where the trade warrants, ranging down to 0.5\%.

The discipline is visible in the executions rather than only in the intention. Of the 16
discretionary positions closed, \textbf{10 sold at or above the pre-set target} and \textbf{6 sold
at the target price itself}, to within two cents --- the signature of a resting limit order filling
at a level set in advance.

\subsection{The complete discretionary record}

\begin{center}
\scriptsize
\begin{tabular}{@{}l l r r r r l r r r@{}}
\toprule
\textbf{Share} & \textbf{Bought} & \textbf{Price} & \textbf{Shares} & \textbf{Cost} &
\textbf{Target} & \textbf{Sold} & \textbf{Price} & \textbf{Hold} & \textbf{P\&L} \\
\midrule
AVGO & 19 Aug & 362.15 & 1,000 & 362,153 & 380.26 & open        & ---    & --- & --- \\
AVGO & 19 Aug & 364.21 & 1,000 & 364,215 & 382.42 & open        & ---    & --- & --- \\
AVGO & 26 Aug & 351.09 & 2,000 & 702,190 & 368.64 & 27 Aug      & 369.20 & 1  & 36,191 \\
MU   & 27 Aug & 912.00 &   500 & 456,002 & 957.60 & 1 Sep       & 967.00 & 5  & 27,493 \\
MU   &  2 Sep & 931.45 &   600 & 558,873 & 978.02 & 2 Sep       & 948.00 & 0  &  9,921 \\
AVGO &  3 Sep & 354.89 & 1,200 & 425,879 & 372.64 & 8 Sep       & 372.00 & 5  & 20,509 \\
VRT  & 10 Sep & 247.50 & 2,000 & 495,010 & 259.88 & 2 Oct       & 252.00 & 22 &  8,971 \\
MU   & 15 Sep & 932.00 &   500 & 466,002 & 978.60 & 17 Sep      & 966.96 & 2  & 17,473 \\
AVGO & 28 Sep & 351.03 & 1,000 & 351,035 & 360.00 & 29 Sep      & 360.00 & 1  &  8,956 \\
VRT  & 28 Sep & 251.64 & 1,300 & 327,139 & 258.64 & open        & ---    & --- & --- \\
VST  & 30 Sep & 135.50 & 3,000 & 406,506 & 142.27 & 2 Oct       & 142.27 & 2  & 20,276 \\
AVGO & 30 Sep & 353.00 & 1,410 & 497,737 & 360.00 & open        & ---    & --- & --- \\
VRT  &  1 Oct & 240.00 & 1,500 & 360,008 & 252.00 & 2 Oct       & 252.00 & 1  & 17,978 \\
VRT  &  1 Oct & 239.51 &   500 & 119,760 & 251.49 & 2 Oct       & 252.00 & 1  &  6,235 \\
VST  &  2 Oct & 139.00 & 3,000 & 417,015 & 139.70 & 2 Oct       & 139.70 & 0  &  2,056 \\
VST  &  2 Oct & 136.00 & 3,000 & 408,015 & 139.70 & 2 Oct       & 139.70 & 0  & 11,056 \\
VST  &  2 Oct & 135.00 &   466 &  62,912 & 139.00 & 2 Oct       & 139.00 & 0  &  1,857 \\
VST  &  2 Oct & 135.00 & 1,034 & 139,595 & 139.00 & 2 Oct       & 139.70 & 0  &  4,845 \\
VST  &  2 Oct & 137.50 & 6,000 & 825,030 & 144.38 & 2 Oct       & 140.00 & 0  & 14,913 \\
VST  &  2 Oct & 137.20 & 1,000 & 137,205 & 144.00 & 2 Oct       & 140.00 & 0  &  2,786 \\
\bottomrule
\end{tabular}
\end{center}

\subsection{The rest of the operation}

\begin{itemize}[leftmargin=1.4em,itemsep=0.2em]
\item \textbf{Ordinary shares only, fully paid.} No derivatives, no options, no futures, no
      contracts for difference, no borrowed stock and no short selling. The company has never held
      a position other than fully paid listed stock.
\item \textbf{A hard maximum holding period of 50 calendar days.} If the exit price has not been
      reached by then, the position is closed at the next opening price regardless of profit or
      loss.
\item \textbf{Two models run on each share}, each allocated half of that share's capital: a
      Bayesian state-space filter and a mean-reversion model. Each takes and closes its own
      positions independently.
\item \textbf{Order placement is automated.} Orders are transmitted to the broker by the company's
      own software each morning before the market opens. Every order carries a machine-readable
      reference identifying the share and the model that generated it, and every resulting
      execution is journalled automatically and reconciled to the broker's records daily.
\item \textbf{Capital is turned over, not accumulated.} Cash released by a sale is redeployed into
      the next purchase on the following morning. Nothing is set aside to be held.
\end{itemize}

\subsection{Funding}

The company is funded by a loan facility from its director, Mr David Fewer, of up to
EUR 10,000,000, of which EUR 5,211,000 was advanced within two months of execution. The facility
is repayable on demand, in whole or in part, and the recital records that the loan is provided
``in order for BC to trade stocks in the stock market and fund its operational costs and
expenses''.

Because the loan is repayable on demand, the company must be able to realise its entire book at
short notice at any time. Its holdings are accordingly confined to large, liquid, exchange-traded
shares and no position may be held longer than 50 days. The financing and the strategy are
consistent with one another.

The loan is interest free. The company's own economic modelling nonetheless charges a cost of
carry of 3.14\% per annum against positions held, so the activity is appraised internally on a
financed basis.

\section{The badges of trade}

Seven of the nine badges support the company's activity being a trade. One is neutral on its face
and is settled by the company's conduct. One does not arise.

\begin{center}
\footnotesize
\begin{tabular}{@{}c l l p{6.4cm}@{}}
\toprule
\textbf{\#} & \textbf{Badge} & \textbf{Indication} & \textbf{Principal facts} \\
\midrule
1 & Profit-seeking motive & Supports & Exit price fixed before entry on 109 of 109 positions;
    purchases only below computed fair value \\
2 & Number and frequency & Supports & 109 positions in 44 trading days across all 9 shares;
    2.48 per trading day \\
3 & Nature of the asset & Neutral & Settled by conduct: no income sought, held in days, turned
    over continuously \\
4 & Expertise; character of operation & Supports & Organised commercial operation: documented
    system, 640 automated checks, daily reconciliation \\
5 & Changes to the asset & Not applicable & No supplementary work is possible on a share \\
6 & Manner of the sale & Supports & Sales executed at a price computed at entry, or forced by the
    50-day rule \\
7 & Source of finance & Supports & Debt funded and repayable on demand; the book must be
    realisable at any time \\
8 & Interval between purchase and sale & Strongly supports & Median 1 day; 38\% same day; longest
    31 days; hard 50-day limit \\
9 & Method of acquisition & Supports & Deliberate purchases at pre-computed limit prices below
    fair value \\
\bottomrule
\end{tabular}
\end{center}

\subsection{Badges 1, 8 and 9 --- intention, holding period and acquisition}

These are one fact seen three ways. A position is created only when the market falls to a price
the company has computed in advance; at the moment it is created its selling price is already
determined; and it cannot survive beyond 50 days.

The realised holding periods bear this out rather than merely asserting it. Of 97 completed
positions, 37 were bought and sold on the same day and 77 were closed within three days. The
longest was held 31 days. Short holding periods are among the most heavily weighted of the badges,
and a median of one day is the clearest possible statement of what the company does with the stock
it buys.

Badge 1 is evidenced by documents rather than by assertion. For the systematic book the sell order
exists in the broker's records within seconds of the buy. For the discretionary book the target is
recorded in the Trade Log before the buy order is transmitted, and in ten of the sixteen closed
positions the stock sold at or above that recorded level.

\subsection{Badge 2 --- number and frequency}

109 positions in the company's first 44 trading days, across all 9 shares in its universe, is a
rate of transaction characteristic of a dealing operation rather than the management of an
investment portfolio. Purchases were placed on 36 of those 44 days. The record is systematic,
fully documented, and the product of a process designed to operate at that frequency indefinitely.

\subsection{Badge 6 --- the manner of the sale}

Sales are not occasioned by a change of view, a need for cash or an opportunity taken. They are
executed at a price computed at entry, or forced by the passage of time under the 50-day rule. The
company's exits are decided before its entries are made.

\subsection{Badge 4 --- expertise and the character of the operation}

The activity is conducted as an organised commercial operation rather than as a series of separate
investment decisions. Every systematic order the company has ever placed was generated by the same
workbook and the same computer code, from the same parameters, and is traceable to it. The company
maintains an operator's manual for the daily process, an automated test suite of approximately 640
checks that must pass before any live run, a trade log reconciled daily against broker records, a
documented research record, and a defined division of responsibilities between three individuals.

\subsection{Badge 3 --- the nature of the asset}

Shares are capable of being held as investments and of being dealt in as stock. The badge is
therefore settled not by the asset class but by what the company does with it, and on that the
record is unambiguous.

\begin{itemize}[leftmargin=1.4em,itemsep=0.2em]
\item \textbf{The company does not seek or rely on the income of the asset.} Dividends are not a
      consideration in selecting shares, form no part of the return the company models, and cannot
      in practice arise in a material way on holdings measured in days.
\item \textbf{The company's return does not derive from appreciation of the asset.} Its own
      analysis, which removes the underlying market trend from the historical record and re-runs
      the system on what remains, finds that approximately 62\% of the return survives the removal
      of that trend --- that is, it arises from short-term price oscillation harvested by repeated
      purchase and resale.
\item \textbf{The holdings are inventory rather than a portfolio.} Cash released by a sale is
      redeployed into the next purchase on the following morning. The same capital is turned over
      repeatedly and nothing is set aside to be held.
\end{itemize}

\subsection{Badge 7 --- source of finance}

The activity is debt funded rather than equity funded, and the debt is repayable on demand. The
company cannot assume it will be able to hold any position to a date of its own choosing. That
constraint is not theoretical --- it is the reason the book consists exclusively of large, liquid
shares and the reason no position may be held beyond 50 days. The financing and the strategy point
the same way.

\section{The operating standard --- what we keep doing}

The company's position rests on what it actually does, recorded as it happens. The following are
the practices that produce that record, and they are to be maintained without exception.

\note{bcact}{\textbf{1. Fix the exit price before entry on every position, in both books.} This is
the single most important practice in the operation. Systematic: the sell order goes to the broker
on execution of the buy. Discretionary: the \emph{Target sale} field is completed in the Trade Log
\emph{before} the buy order is transmitted --- not after the fill is confirmed. The sequence is
what matters. Retain the 5.0\% default and continue recording tighter targets where the trade
warrants one.}

\note{bcact}{\textbf{2. Never make an exception to the 50-day maximum hold.} The rule's value lies
entirely in its being absolute. The longest completed hold to date is 31 days, so the rule has
never yet cost the company anything --- which is precisely why it must be honoured the first time
it does. Review any position passing 35 days.}

\note{bcact}{\textbf{3. Keep the book active.} The current rate is 2.48 new positions per trading
day, with purchases on 36 of 44 trading days and dealings in all nine shares. A quiet book weakens
the clearest badge the company has. Where the models generate no orders for an extended period, the
reason should be recorded at the time.}

\begin{itemize}[leftmargin=1.4em,itemsep=0.25em]
\item \textbf{Deal only in fully paid listed ordinary shares.} No derivatives, options, futures,
      contracts for difference, borrowed stock or short positions. A single instrument outside this
      description changes the character of the activity.
\item \textbf{Buy below the market, always.} Limit orders at pre-computed prices. No market orders
      and no raising the bid to obtain stock.
\item \textbf{Redeploy released capital the following morning.} The capital is inventory funding,
      not a portfolio allocation.
\item \textbf{Reconcile the trade log to broker execution records daily}, and keep the
      machine-readable order references on every order.
\item \textbf{Run the automated test suite before every live session.} All 640 checks must pass.
\item \textbf{Keep dividends out of share selection.} They form no part of the modelled return and
      should form no part of the decision.
\item \textbf{Keep every order traceable} to the workbook, the code and the parameters that
      generated it.
\item \textbf{Describe the activity consistently} across the website, the financial statements,
      board minutes and internal documents.
\item \textbf{Maintain the research record.} The testing and development history evidences an
      organised commercial operation.
\end{itemize}

\section{Maintaining the record}

This document is rewritten each month from the Trade Log and the Discretionary Trade Log in the
current trading workbook. The monthly refresh covers:

\begin{itemize}[leftmargin=1.4em,itemsep=0.15em]
\item Positions opened, closed and open; same-day and three-day counts; median and longest holding
      period; losses; gross purchases; largest position; realised profit.
\item Confirmation that every position opened in the month carried an exit price fixed before
      entry, in both books.
\item Confirmation that no position exceeded the 50-day limit.
\item The weekly realised profit table and the positions-by-share table.
\item The date of the most recent reconciliation to broker records.
\end{itemize}

\subsection{Revision history}

\begin{center}
\footnotesize
\begin{tabular}{@{}l l r p{7.2cm}@{}}
\toprule
\textbf{Date} & \textbf{Record to} & \textbf{Positions} & \textbf{Changes} \\
\midrule
5 Oct 2026 & 2 Oct 2026 & 109 & Rewritten as the company's internal record. Discretionary book
  confirmed to carry an exit price fixed before entry on all 20 positions. Badge 1 stated without
  qualification. Operating standard added. \\
28 Sep 2026 & 28 Sep 2026 & 92 & Discretionary book brought into the document. \\
27 Sep 2026 & 25 Sep 2026 & 80 & First version. \\
\bottomrule
\end{tabular}
\end{center}

\vspace{0.6em}
{\color{bcrule}\rule{\textwidth}{0.6pt}}

{\footnotesize\color{bcgrey}
Trading figures are taken from the Trade Log and Discretionary Trade Log in
\texttt{TradingExcel\_9stock\_avgo.xlsx} as at 2 October 2026, covering 4 August to 2 October 2026
and reconciled to broker execution records. Loan terms are taken from the executed loan agreement
between Mr David Fewer and the company. The trend-removal analysis referred to at badge 3 is
recorded in the company's research file.}

\end{document}
